\begin{theorem}\label{thm:kernel-chebotarev-quotient-relative-entropy-chain}
Let \(\pi:G_2\twoheadrightarrow G_1\) be a surjective homomorphism of finite abelian groups, and write \(N:=\ker\pi\).
Fix \(u\in(0,1]\) and \(n\ge 1\), and suppose that
\[
B_n(\ind;u):=
\sum_{c_2\in G_2}B^{(2)}_{n,c_2}(u)
=
\sum_{c_1\in G_1}B^{(1)}_{n,c_1}(u)
>0,
\]
where the coarse counts are obtained from the fine counts by pushforward:
\[
B^{(1)}_{n,c_1}(u):=
\sum_{c_2\in\pi^{-1}(c_1)}B^{(2)}_{n,c_2}(u)
\qquad(c_1\in G_1).
\]
Define probability measures
\[
\mu_n^{(i)}(c):=\frac{B^{(i)}_{n,c}(u)}{B_n(\ind;u)}
\qquad(i=1,2),
\]
and let \(\mathrm{unif}_{G_i}\) denote the uniform measure on \(G_i\). For each \(c_1\in G_1\) with \(\mu_n^{(1)}(c_1)>0\), define
\[
\mu_{n,\mathrm{fib}}^{(c_1)}(c_2):=
\frac{\mu_n^{(2)}(c_2)}{\mu_n^{(1)}(c_1)}
\qquad(c_2\in\pi^{-1}(c_1)),
\]
and write \(\mathrm{unif}_{\pi^{-1}(c_1)}\) for the uniform measure on the fiber \(\pi^{-1}(c_1)\). Then
\[
D\!\bigl(\mu_n^{(2)}\,\|\,\mathrm{unif}_{G_2}\bigr)
=
D\!\bigl(\mu_n^{(1)}\,\|\,\mathrm{unif}_{G_1}\bigr)
+
\sum_{c_1\in G_1}\mu_n^{(1)}(c_1)\,
D\!\bigl(\mu_{n,\mathrm{fib}}^{(c_1)}\,\|\,\mathrm{unif}_{\pi^{-1}(c_1)}\bigr).
\]
As usual, fibers with \(\mu_n^{(1)}(c_1)=0\) contribute \(0\).
\end{theorem}

\begin{proof}
This is the chain rule for a partition into fibers, once one notes that the reference measure factors accordingly. By construction,
\[
\mu_n^{(1)}=\pi_\ast\mu_n^{(2)}.
\]
Since \(\pi\) is a surjective homomorphism of finite groups, every fiber has cardinality \(|N|\), and therefore
\[
\mathrm{unif}_{G_2}(c_2)
=
\frac{1}{|G_2|}
=
\frac{1}{|G_1|\,|N|}
=
\mathrm{unif}_{G_1}(\pi(c_2))\,
\mathrm{unif}_{\pi^{-1}(\pi(c_2))}(c_2).
\]
If \(\mu_n^{(1)}(c_1)=0\), then every term \(\mu_n^{(2)}(c_2)\) with \(c_2\in\pi^{-1}(c_1)\) is zero, since their sum is \(\mu_n^{(1)}(c_1)\). Such fibers therefore contribute nothing.

For \(c_2\in G_2\) with \(\mu_n^{(1)}(\pi(c_2))>0\), the fine distribution factors as
\[
\mu_n^{(2)}(c_2)
=
\mu_n^{(1)}(\pi(c_2))\,
\mu_{n,\mathrm{fib}}^{(\pi(c_2))}(c_2).
\]
Hence
\[
\log\frac{\mu_n^{(2)}(c_2)}{\mathrm{unif}_{G_2}(c_2)}
=
\log\frac{\mu_n^{(1)}(\pi(c_2))}{\mathrm{unif}_{G_1}(\pi(c_2))}
+
\log\frac{\mu_{n,\mathrm{fib}}^{(\pi(c_2))}(c_2)}{\mathrm{unif}_{\pi^{-1}(\pi(c_2))}(c_2)}.
\]
Multiply by \(\mu_n^{(2)}(c_2)\), sum over \(c_2\in G_2\), and regroup by fibers of \(\pi\). This yields the stated identity.
\end{proof}

\begin{theorem}\label{thm:kernel-chebotarev-visible-hidden-character-energy}
In the setting of Theorem \ref{thm:kernel-chebotarev-quotient-relative-entropy-chain}, define the relative densities
\[
r_n^{(i)}(c):=|G_i|\,\mu_n^{(i)}(c)
=
\frac{|G_i|\,B^{(i)}_{n,c}(u)}{B_n(\ind;u)}
\qquad(i=1,2).
\]
Then
\[
\|r_n^{(2)}-1\|_{L^2(\mathrm{unif}_{G_2})}^2
=
\|r_n^{(1)}-1\|_{L^2(\mathrm{unif}_{G_1})}^2
+
\frac{1}{|G_2|}\sum_{c_1\in G_1}\sum_{c_2\in\pi^{-1}(c_1)}
\bigl(r_n^{(2)}(c_2)-r_n^{(1)}(c_1)\bigr)^2.
\]
If
\[
N^\perp:=\{\chi\in\widehat G_2:\chi|_N\equiv 1\}\cong \widehat G_1
\]
and
\[
B_n^{(2)}(\chi;u):=\sum_{c_2\in G_2}\chi(c_2)\,B^{(2)}_{n,c_2}(u)
\qquad(\chi\in\widehat G_2),
\]
then
\[
\|r_n^{(2)}-1\|_{L^2(\mathrm{unif}_{G_2})}^2
=
\frac{1}{B_n(\ind;u)^2}
\sum_{\chi\in\widehat G_2\setminus\{\mathbf{1}\}}
\bigl|B_n^{(2)}(\chi;u)\bigr|^2,
\]
\[
\|r_n^{(1)}-1\|_{L^2(\mathrm{unif}_{G_1})}^2
=
\frac{1}{B_n(\ind;u)^2}
\sum_{\chi\in N^\perp\setminus\{\mathbf{1}\}}
\bigl|B_n^{(2)}(\chi;u)\bigr|^2.
\]
Consequently,
\[
\|r_n^{(2)}-1\|_{L^2(\mathrm{unif}_{G_2})}^2
\;-\;
\|r_n^{(1)}-1\|_{L^2(\mathrm{unif}_{G_1})}^2
\;=\;
\frac{1}{B_n(\ind;u)^2}
\sum_{\chi\in\widehat G_2\setminus N^\perp}
\bigl|B_n^{(2)}(\chi;u)\bigr|^2.
\]
\end{theorem}

\begin{proof}
The relevant observation is that fiberwise averaging is the orthogonal projection onto the subspace of functions that are constant on the fibers of \(\pi\). For \(c_1\in G_1\), the pushforward relation reads
\[
r_n^{(1)}(c_1)
=
\frac{1}{|N|}\sum_{c_2\in\pi^{-1}(c_1)}r_n^{(2)}(c_2).
\]
Let \(\mathcal H_\pi\subset L^2(\mathrm{unif}_{G_2})\) be the subspace of functions constant on the fibers of \(\pi\). The displayed identity shows that \(r_n^{(1)}\circ\pi\in\mathcal H_\pi\). If \(h\in\mathcal H_\pi\) and \(h(c_2)=h_0(\pi(c_2))\), then
\[
\begin{aligned}
\left\langle r_n^{(2)}-r_n^{(1)}\circ\pi,h\right\rangle_{L^2(\mathrm{unif}_{G_2})}
&=
\frac{1}{|G_2|}\sum_{c_1\in G_1}\sum_{c_2\in\pi^{-1}(c_1)}
\bigl(r_n^{(2)}(c_2)-r_n^{(1)}(c_1)\bigr)\overline{h_0(c_1)}\\
&=
\frac{1}{|G_2|}\sum_{c_1\in G_1}
\left(\sum_{c_2\in\pi^{-1}(c_1)}r_n^{(2)}(c_2)-|N|\,r_n^{(1)}(c_1)\right)
\overline{h_0(c_1)}\\
&=0.
\end{aligned}
\]
Thus \(r_n^{(1)}\circ\pi\) is the orthogonal projection of \(r_n^{(2)}\) onto \(\mathcal H_\pi\). The Pythagorean theorem in \(L^2(\mathrm{unif}_{G_2})\) gives
\[
\|r_n^{(2)}-1\|_{L^2(\mathrm{unif}_{G_2})}^2
=
\|r_n^{(1)}\circ\pi-1\|_{L^2(\mathrm{unif}_{G_2})}^2
+
\|r_n^{(2)}-r_n^{(1)}\circ\pi\|_{L^2(\mathrm{unif}_{G_2})}^2.
\]
Because each fiber has cardinality \(|N|\), the first term is exactly
\[
\|r_n^{(1)}-1\|_{L^2(\mathrm{unif}_{G_1})}^2,
\]
and expanding the second term fiberwise yields the first displayed identity.

For the Fourier identities, apply Parseval on \(G_2\). The Fourier coefficient of \(r_n^{(2)}\) at \(\chi\in\widehat G_2\) is
\[
\widehat r_n^{(2)}(\chi)
:=
\frac{1}{|G_2|}\sum_{c_2\in G_2}r_n^{(2)}(c_2)\,\overline{\chi(c_2)}
=
\frac{B_n^{(2)}(\overline\chi;u)}{B_n(\ind;u)}.
\]
Since complex conjugation permutes \(\widehat G_2\), Parseval gives
\[
\|r_n^{(2)}-1\|_{L^2(\mathrm{unif}_{G_2})}^2
=
\sum_{\chi\in\widehat G_2\setminus\{\mathbf{1}\}}\abs{\widehat r_n^{(2)}(\chi)}^2
=
\frac{1}{B_n(\ind;u)^2}
\sum_{\chi\in\widehat G_2\setminus\{\mathbf{1}\}}
\bigl|B_n^{(2)}(\chi;u)\bigr|^2.
\]
Likewise, the characters of \(G_1\) identify under pullback with \(N^\perp\subseteq \widehat G_2\), and the same argument gives
\[
\|r_n^{(1)}-1\|_{L^2(\mathrm{unif}_{G_1})}^2
=
\frac{1}{B_n(\ind;u)^2}
\sum_{\chi\in N^\perp\setminus\{\mathbf{1}\}}
\bigl|B_n^{(2)}(\chi;u)\bigr|^2.
\]
Subtracting the two identities yields the hidden-energy formula.
\end{proof}

\begin{remark}
Nothing here is new at the level of information theory or harmonic analysis; compare \cite{CoverThomas2006,CsiszarKorner2011,Diaconis1988,Terras1999}. The point in the present setting is that the coarse counts arise by exact pushforward from the primitive layer, so the discarded term acquires a canonical arithmetic meaning: it is precisely the character energy outside \(N^\perp\).
\end{remark}

\endinput
